% =====================================================================
% sections/01-introduction.tex  --  Member A
% Labels used by other sections: sec:introduction
% Cross-references to other members' sections (labels agreed by team):
%   sec:mlops-concepts, sec:michelangelo-overview (A)
%   sec:design-rationale (B), sec:pros-cons (C),
%   sec:layered-architecture (D), sec:conclusion (A)
% =====================================================================

\section{Introduction}
\label{sec:introduction}

Machine learning (ML) has moved from a research activity to a core
component of production software, yet the engineering effort required to
operate ML systems is routinely underestimated. Sculley et al.\ observe
that developing and deploying ML systems is relatively fast and cheap,
whereas maintaining them over time is difficult and expensive, because
ML systems inherit all the maintenance problems of conventional code and
add a set of ML-specific ones, such as data dependencies, feedback loops
and sensitivity to changes in the external world~\cite{sculley2015hidden}.
In a typical real-world ML system, the model code is only a small part;
most of the system consists of surrounding infrastructure for data
collection, feature extraction, serving and monitoring~\cite{sculley2015hidden}.
Empirical studies confirm that the ML workflow is data-centric and highly
non-linear, with frequent feedback loops between data, modelling and
deployment stages~\cite{amershi2019software}. \emph{Machine Learning
Operations} (MLOps) has emerged as the engineering paradigm that addresses
this gap by combining ML, software engineering (in particular DevOps) and
data engineering to bring ML products into production and keep them
there~\cite{kreuzberger2023mlops}.

This report studies MLOps through the lens of \emph{software architecture},
using Uber's internal ML platform \emph{Michelangelo} as a case study.
Before Michelangelo, data scientists at Uber built models with a
heterogeneous set of tools, while engineering teams built bespoke,
one-off systems to run those models in production; there was no
standard way to store and compare experiments and no established path
to deployment~\cite{delbalso2017meet}. Michelangelo was created to
standardise this end-to-end workflow~\cite{delbalso2017meet} and has
since evolved through three phases: predictive ML on tabular data,
deep learning (DL), and generative AI~\cite{wang2024predictive}. We
select Michelangelo for three reasons. First, it is one of the earliest
documented enterprise ML platforms, and its feature store is described
by authors involved in building it as the first industrial feature
store~\cite{orr2021managing}. Second, its design and evolution are
documented both in peer-reviewed publications~\cite{li2017scaling,sun2020gallery}
and in a series of engineering reports spanning eight
years~\cite{delbalso2017meet,delbalso2018scaling,wang2024predictive},
including Uber's own account of mistakes and lessons learned. Third, it
operates at a scale that makes architectural trade-offs visible: as of
2024, the platform manages about 400 active ML projects, runs more than
20{,}000 training jobs per month, and hosts over 5{,}000 production models
that serve up to 10 million real-time predictions per second at
peak~\cite{wang2024predictive}.

The report addresses four questions: (i)~which MLOps concepts a platform
such as Michelangelo must realise; (ii)~why Michelangelo was designed the
way it is; (iii)~what trade-offs this design entails; and (iv)~how its
layered architecture organises the platform's responsibilities.
Section~\ref{sec:mlops-concepts} defines the core MLOps concepts and maps
each to its Michelangelo counterpart, and
Section~\ref{sec:michelangelo-overview} summarises the platform's evolution
from 2015 to 2024 together with its six-step workflow.
Section~\ref{sec:design-rationale} derives the main design decisions from
the problems Uber faced in 2015. Section~\ref{sec:pros-cons} discusses
the platform's strengths, weaknesses and trade-offs, including a brief
comparison with alternative platforms. Section~\ref{sec:layered-architecture}
analyses Michelangelo as a layered architecture, and
Section~\ref{sec:conclusion} concludes the report.
% TODO (A): adjust the sentence above if B/C/D change the scope of their sections.
% TODO (A): if G's demo is described in the report, add one sentence here
%           and a \ref to the corresponding section/appendix.